\chapter{Glaucoma Tests}
\section*{What Is This Test?} Glaucoma testing assesses the optic nerve, visual field, eye pressure, cornea, and drainage angle for damage or risk.\section*{Alternative Names} Glaucoma screening; tonometry; optic-nerve evaluation.\section*{Principle} Tonometry measures intraocular pressure, while ophthalmoscopy, optical coherence tomography, perimetry, and gonioscopy assess structure and function.\section*{Why This Test Is Ordered} It evaluates raised eye pressure, optic-nerve appearance, vision loss, family history, or suspected glaucoma.\section*{Primary vs Secondary Test} Eye examination is primary; pressure alone cannot diagnose or exclude glaucoma.\section*{How This Test Is Performed} The clinician measures pressure, examines the optic nerve, maps visual fields, and may image retinal nerve fibers or inspect the drainage angle.\section*{What Preparation Is Needed} Contact lenses may need removal. Bring glasses, medicines, and prior eye records.\section*{Understanding Results} Glaucoma diagnosis depends on structure, function, pressure, and progression; normal pressure does not exclude disease.\section*{Follow-up Tests} Repeat fields, OCT, gonioscopy, pachymetry, or ophthalmology follow-up may be needed.\section*{References} \begin{enumerate}\item MedlinePlus. Glaucoma Tests.\item American Academy of Ophthalmology. Primary open-angle glaucoma guidance.\end{enumerate}\clearpage\section*{Understanding Results and Follow-up}
The laboratory report should identify the specimen, method, units, reference interval or decision limit, and whether the result is positive, negative, abnormal, or inconclusive. Reference intervals vary with age, sex, pregnancy, specimen type, assay, and laboratory; a result outside a range is not automatically a diagnosis, and a result within a range does not always exclude disease. Discuss unexpected results with the ordering clinician, who can decide whether repeat or confirmatory testing, treatment, or referral is appropriate. Do not change medicines or delay urgent care based on an isolated result.
\section*{Regulatory Status and Authorization Date}This chapter describes a test category, not a specific commercial product. FDA, CDSCO, EU IVDR, and MHRA approval or registration dates therefore cannot be assigned without the assay manufacturer, product name, instrument, and intended use. For laboratory-developed tests, an individual agency approval date may not exist. See the Preface for the interpretation of regulatory status.
\clearpage
